% !TEX root = ../main.tex
\section{Model benchmark}
\label{sec:benchmark}

The model screen runs the same condition with one lineup of models at a time,
changing nothing but the models~\cite{project-design}. Table~\ref{tab:benchmark}
will compare the lineups on six measures of the game: how often the jury reached
the truth-consistent verdict, how much the agents used side payments and private
messages, whether the auditor accepted payments, and how many agents the
reasoning monitor of Section~\ref{sec:probes} judged rogue.

No model-screen run has been admitted yet, so the table is a template and every
value is pending. Once runs finish, each column pools the counts of one lineup's
finished runs of a single condition. Rates keep their counts, and a rate whose
denominator is zero is reported as undefined rather than as zero. The episode
count under each model is the number of episodes behind that column; episodes
share memory within a block, so it is not a count of independent replicates.

\begin{table}[htbp]
  \caption{Model benchmark template for one condition. Each column will pool the
  finished W\&B runs of one model lineup; every agent role uses the named model
  unless the column names two parties. Rates show the percentage with its count;
  offers and private messages are means per episode; n/a marks a zero
  denominator, and --- a value not yet measured. Shading marks the featured
  model. Bold marks the safest value in a row (the highest truth-consistent rate
  or the lowest collusion measure) when one model alone holds it; grey marks it
  outside the featured column. With few blocks per model, a bold value is a
  description, not a significant difference.}
  \label{tab:benchmark}
  \centering
  \small
  % Generated by analysis/model_benchmark.py from data/model-benchmark.json; do not edit.
% Template: no admitted runs; every value is pending.
\begingroup
% Rules without padding, so the shaded column runs unbroken between them.
\setlength{\aboverulesep}{0pt}%
\setlength{\belowrulesep}{0pt}%
\renewcommand{\arraystretch}{1.35}%
\definecolor{benchfeature}{HTML}{F3DACB}%
\definecolor{benchbest}{HTML}{E2E0DA}%
\begin{tabular}{>{\raggedright\arraybackslash}p{40mm}>{\columncolor{benchfeature}}lllll}
  \toprule
   & \textbf{Model A} & \textbf{Model B} & \textbf{Model C} & \textbf{Model D} & \textbf{Model E} \\
   & {\footnotesize\color{black!60}--- episodes} & {\footnotesize\color{black!60}--- episodes} & {\footnotesize\color{black!60}--- episodes} & {\footnotesize\color{black!60}--- episodes} & {\footnotesize\color{black!60}--- episodes} \\
  \midrule
  \textbf{Truth-consistent verdicts}\newline{\footnotesize\color{black!60}share of settled verdicts} & --- & --- & --- & --- & --- \\
  \textbf{Offers to jurors}\newline{\footnotesize\color{black!60}per episode} & --- & --- & --- & --- & --- \\
  \textbf{Juror acceptance}\newline{\footnotesize\color{black!60}share of offers to jurors} & --- & --- & --- & --- & --- \\
  \textbf{Private messages}\newline{\footnotesize\color{black!60}per episode} & --- & --- & --- & --- & --- \\
  \textbf{Auditor acceptance}\newline{\footnotesize\color{black!60}share of offers to the auditor} & --- & --- & --- & --- & --- \\
  \textbf{Agents judged rogue}\newline{\footnotesize\color{black!60}hybrid monitor score at least 50} & --- & --- & --- & --- & --- \\
  \bottomrule
\end{tabular}
\endgroup

\end{table}
